\documentclass[10pt,a4paper]{amsart}
\usepackage[margin=19mm]{geometry}
\usepackage{amsmath,amssymb,mathtools}
\usepackage[scr=rsfs,cal=euler]{mathalfa}
\usepackage{xfrac}
\usepackage[unicode,hidelinks,bookmarks=false]{hyperref}
\newcommand{\ie}{i.e.\ }
\newcommand{\cf}{cf.\ }
\newcommand{\resp}{respectively\ }
\newcommand{\Subsection}{\subsection}
\providecommand{\nobreakdash}{\nobreak\hbox{-}}
\setlength{\emergencystretch}{2em}
\multlinegap0pt
\usepackage{tikz}
\usetikzlibrary{cd}
\usepackage{amssymb}
\usepackage{tikz}
\usepackage{enumerate}
\usepackage[shortlabels]{enumitem}
\usepackage{mathtools}
\usepackage{float}
\newtheorem{theorem}{Theorem}
\newtheorem{lemma}{Lemma}
\newcommand{\CC}{\cC \cC}
\newcommand{\IC}{\cI \cC}
\newcommand{\cC}{\mathcal C}
\newcommand{\cI}{\mathcal I}
\newcommand{\cN}{\mathcal N}
\newcommand{\ga}{\alpha}
\title{Irreducible Characteristic Cycles for Orbit Closures of a Symmetric Subgroup}
\author{William Graham, Minyoung Jeon, Scott Joseph Larson}
\date{Source: arXiv:2603.28392v1 (CC BY 4.0)}
\begin{document}
\maketitle

\noindent\emph{Source and licence.} Irreducible Characteristic Cycles for Orbit Closures of a Symmetric Subgroup. Version arXiv:2603.28392v1 (CC BY 4.0), distributed by arXiv under the Creative Commons Attribution 4.0 International licence, http://creativecommons.org/licenses/by/4.0/, as recorded in the arXiv licence metadata for this exact version and shown on the abstract page https://arxiv.org/abs/2603.28392v1. Full licence notice: \texttt{licenses/arxiv-2603.28392-CC-BY-4.0.txt}.\par\smallskip

\noindent\textit{Editorial scope.} Counted unit: the article's theorem t:symmetric-mu (source0.tex lines 1473--1494). The article's complete original proof of it is delivered with it (source0.tex lines 1609--1665, one \texttt{proof} environment of 57 lines, which the article places away from the statement). No mathematics is written, repaired or re-derived: the statement and the proof are the article's own lines, in the article's own notation, and no retained line is substituted or re-flowed.  Retained uncounted as the source-local context the proof needs: lines 1257--1259; lines 1423--1427; lines 1540--1547; lines 1599--1602.  Nothing was omitted from any retained span, so the package carries no omission record.  No retained line contains a citation, so the wrapper carries no bibliography and no bibliography entry.  The retained lines contain the article's own diagram code, which the wrapper typesets with the standard tikz packages; no image file is delivered and the package declares no asset dependency.  Editorial formatting only, in addition: an AMS \texttt{amsart} preamble that restates the article's own declarations and macros used by the retained lines, namely the environments \texttt{theorem}, \texttt{lemma} on separate counters, together with the standard packages those lines need.  The retained environments print in this document as Theorem 1, Lemma 1 in order of appearance, whereas the article prints the same items with the numbering of its own full document.  The complete native source of the article as distributed in the arXiv e-print for this exact version is bundled unchanged as \texttt{sources/197432/source0.tex}.
\begin{equation} \label{e:badconfig}
\epsilon | a a, \hspace{.2in} a a | \epsilon,  \hspace{.2in} \epsilon | a \epsilon,  \hspace{.2in} \epsilon a | \epsilon.
\end{equation}

\begin{theorem} \label{t:irreducible}
Let $G = GL(n)$ and $K = GL(p) \times GL(q)$.  Let $X_{v_0}$ be a smooth orbit closure in $G/B$, and let $I$ be a set of commuting simple reflections such that 
$\ell(v_0* w_I) = \ell(v_0)+ |I|$ and $\ell(y* w_I) = \ell(y) + |I|$ for all $y \leq v_0$ with $\ell(y) = \ell(v_0) - 1$.  Then for $v=v_0*w_I$, the characteristic
cycle $\CC(\IC_{X_v})$ is irreducible.
\end{theorem}

\begin{lemma} \label{l:tau}
Assume that the hypotheses of Theorem \ref{t:irreducible}  hold.  Then 
every element of $\tau(v)$
commutes with every element of $I$, and
\begin{equation} \label{e:tau}
\tau(v) = (\tau(v) \cap \tau(v_0)) \sqcup I.
\end{equation}
\end{lemma}

\begin{lemma} \label{l:fiber-small}
Let $\mu_1: Z_1 \to Y$ and $\mu_2: Z_2 \to Y$ be small resolutions of $Y$.  Then for any $y \in Y$,
$H^*(\pi_1^{-1}(y)) = H^*(\pi_2^{-1}(y))$.  Hence $\dim \pi_1^{-1}(y) = \dim \pi_2^{-1}(y)$.
\end{lemma}

\begin{theorem}\label{t:symmetric-mu}
Assume that the hypotheses of Theorem \ref{t:irreducible}  hold.
Let $\cN$ be a subset of $I$ consisting of $v_0$-noncompact imaginary simple reflections, and
let $I' = I \smallsetminus \cN$.  
Let $v_0' = v_0 * w_{\cN}$, and \(M= \tau(v_0') \cap\tau(v)\).
There is a commutative diagram
\begin{equation}\label{e:symmetric-mu}
\begin{tikzcd}
Z_{v_0, I} = G_{v_0} \times^B P_I/B\arrow[r,"\mu"]\arrow[d,swap,"\varphi"]&\overline{KvB}/B\\
G_{v_0'} \times^{P_M} P_{\tau(v)}/B\arrow[ur,swap,"\mu'"]
\end{tikzcd}
\end{equation}
such that \(\varphi\) is a \(K\)-equivariant isomorphism.  In particular, taking $\cN$ to be empty,
we obtain an isomorphism
\begin{equation}\label{e:symmetric-mu2}
Z_{v_0, I} \cong G_{v_0} \times^{P_M} P_{\tau(v)}/B.
\end{equation}
Moreover, for general $\cN$, we have
\begin{equation}\label{e:symmetric-mu3}
Z_{v_0, I} \cong Z_{v_0', I'}.
\end{equation}
\end{theorem}

\begin{proof}[Proof of Theorem \ref{t:symmetric-mu}]
We first show that $\varphi$ is an isomorphism under the assumption that $\cN$ is empty.  
Since $G_{v_0} = K \dot{v_0} P_{\tau(v_0)}$  is right-invariant under $P_M \subset P_{\tau(v_0)}$,
and $I \subset \tau(v)$, there is a morphism $\varphi: Z = G_{v_0} \times^B P_I/B \to G_{v_0} \times^{P_M} P_{\tau(v)}/B$
given by \(\varphi[g,p]=[g,p]\).  This morphism is proper since both source and target are complete.
We claim that to show that $\varphi$ is an isomorphism, it is enough to show that $\varphi$ is surjective.
Indeed, suppose $\varphi$ is surjective.  Then $\mu'$ is also a small resolution. By Lemma \ref{l:fiber-small},
over any given point in $X_v$, the fiber of $\mu$ and the fiber of $\mu'$ have the same dimension, 
so $\varphi$ is a finite morphism.
Since $\varphi$ is $K$-equivariant, it must be a covering
over the open $K$-orbit $K \dot{v} B/B$ in $X_v$.  Since these $K$-orbits are equivariantly simply connected, this implies $\varphi$ is an isomorphism over
this $K$-orbit, so $\varphi$ is birational.  Since both source and target of $\varphi$ are smooth, hence
normal, Zariski's connectedness theorem implies that the fibers of $\varphi$ are connected; then Zariski's Main Theorem implies that
$\varphi$ is an isomorphism.  This proves the claim.

Since $\mu = \mu' \circ \varphi$ and $\mu$ is birational, we have $\dim \varphi(Z) = \dim Z$.  On the
other hand,
\begin{equation} \label{e:dimension}
\dim K \dot{v}_0P_{\tau(v_0)}\times^{P_M} P_{\tau(v)}/B = \ell(v_0) + \ell(w_{\tau(v)}) - \ell(w_M) = \ell(v_0) + | I | = \dim Z.
\end{equation}
where the last equality follows from Lemma \ref{l:tau}.  Hence $\varphi$ is dominant; since $\varphi$ is proper
and $K\dot{v}_0P_{\tau(v_0)}\times^{P_M} P_{\tau(v)}/B$ is irreducible, we conclude that $\varphi$ is surjective.
This proves that $\varphi$ is an isomorphism in case $\cN$ is empty; in particular, it proves that  \eqref{e:symmetric-mu2} is an isomorphism.

We now consider the general case.  By induction on $|\cN|$, it suffices to prove the result in case $\cN = \{ s_{i} \}$,
so $v_0' = v_0 * w_{\cN} = v_0 * s_{i}$.  As noted before the proof, the only positions where
the clan of $v_0$ differs from that of $v_0'$ are in positions $i, i+1$; in $v_0$, the entries are opposite signs, and
in $v_0'$ (and in $v$), they are equal numbers.  The clan $v_0'$ has a block decomposition where
the blocks are of the form $\max(V_{p_j,q_j})$ (since $v_0$ does), so $X_{v_0'}$ is smooth.  
Since the pair $v_0, I$ avoids the four configurations of  \eqref{e:badconfig}, so does the pair $v_0', I' = I \smallsetminus \cN$,
so this pair satisfies the hypotheses of Theorem \ref{t:irreducible}.  Since $v_0\leq v_0'$,
we have $K \dot{v}_0 P_{\tau(v_0)} = \overline{K \dot{v}_0B} \subset\overline{K{v_{0}'}B} = K v_0' P_{\tau(v_0')}$.

We claim that $\tau(v_0) \cap \tau(v) \subseteq \tau(v_0') \cap \tau(v)$. Indeed, the only positions where
the clan of $v_0$ differs from that of $v_0'$ are in positions $i, i+1$, so
the only possible roots which are in $\tau(v_0)$ but not in $\tau(v_0')$
are $\ga_{i-1}$ and $\ga_{i+1}$.  Neither of these roots is in $\tau(v)$; the claim follows.

Let $M_0 = \tau(v_0) \cap \tau(v)$ and $M = \tau(v_0') \cap \tau(v)$.  We have
$$
\varphi:  Z_{v_0, I} \cong K \dot{v}_0 P_{\tau(v_0)} \times^{P_{M_0}} P_{\tau(v)}/B \to K \dot{v}_0' P_{\tau(v_0')} \times^{P_{M}} P_{\tau(v)}/B \cong Z_{v_0', I'},
$$
where the two congruences follow from the first part of the proof, and the existence of the map
in the middle follows from the discussion above.
Both $v_0, I$ and $v_0', I'$ satisfy the hypotheses of Theorem \ref{t:irreducible}, so by the first part of the proof,
both source and target of $\varphi$ are small resolutions of $X_v$.  Moreover, $\varphi$
commutes with the maps to $X_v$.   Both source and target of $\varphi$
are smooth, complete and irreducible, and the maps to $X_v$ are birational.
Hence, reasoning as in the first part of the proof
shows that $\varphi$ is an isomorphism.  

Finally, we prove that \eqref{e:symmetric-mu3} is an isomorphism.  Applying equation \eqref{e:symmetric-mu2}
to $v_0', I'$ implies
that the target of $\varphi$ in Theorem \ref{t:symmetric-mu}
is isomorphic to $Z_{v_0', I'}$.  Hence $\varphi$ gives an isomorphism
$Z_{v_0,I} \to Z_{v_0', I'}$, as desired. 
\end{proof}

\end{document}
